Recent work has established that mechanistic interpretability approaches can be used to interpret and steer diffusion models during inference. In this work, we train Sparse Autoencoders (SAEs) on the diffusion module and token initializer of RFDiffusion3, a popular all-atom diffusion model for the design of de novo proteins. We identify features for virulence, toxicity, and immunogenicity, using these features for downstream tasks like classification and steerable protein design. Overall, we show RFDiffusion3 has learned useful representations for both token-level and atom-level features, unlocking the internals of the model for further study and controllable sampling. \end{abstract}

\section{Introduction}
Deep learning (DL) models have seen increasing capabilities and popularity in the field of computational biology, especially in the prediction and design of protein structure. Models such as the Alphafold and RosettaFold series excel at predicting protein structure from amino acid sequences, while protein design models such as RFDiffusion, Protenix, and DISCO generate new protein structures conforming to input constraints. Generative protein design models often contain diffusion transformers, inspired by approaches from generative image models \cite{NEURIPS2020_4c5bcfec}. Typically, model outputs are guided by conditioning signals during training, for example reference protein fragments (motifs) and hydrogen donor/acceptor status of atoms. Adherence to these constraints can be amplified with techniques such as classifier-free guidance \cite{ho2021classifierfree}. However, the uninterpretable nature of these models make inference time guidance on novel conditions difficult. Another limitation is the large number of samples generated and filtering required to find effective designs.

Concurrently, the field of mechanistic interpretability aims to deconstruct the intermediate calculations or activations of neural networks into human understandable features, and use these features to explain model outputs. Approaches include training sparse transformers, classifiers on activations, cross and transcoders, and sparse autoencoders. A large body of work has been established around Large Language Model (LLM) interpretability, with recent works focusing on protein language models (PLMs) and scientific models. 

In this work, we train sparse autoencoders on RFDiffusion3. Our goal is to find interpretable SAE features in RFDiffusion3 that describe meaningful structural and functional attributes of proteins. These features can be used as tunable hyperparameters during inference, pioneering a new paradigm for interpretable and explainable protein design. 